Reward models have become a staple in modern NLP, serving as not only a scalable text evaluator, but also an indispensable component in many alignment recipes and inference-time algorithms.
However, while recent reward models increase performance on standard benchmarks, this may partly be due to overfitting effects, which would confound an understanding of their true capability.
In this work, we scrutinize the robustness of reward models and the extent of such overfitting.
We build \textbf{\benchmark}, which systematically transforms reward model inputs in meaning- or ranking-preserving ways.
We show that state-of-the-art reward models suffer from substantial performance degradation even with minor input transformations, sometimes dropping to significantly below-random accuracy, suggesting brittleness.
To improve reward model robustness, we propose to explicitly train them to assign similar scores to paraphrases, and find that this approach also improves robustness to other distinct kinds of transformations.
For example, our robust reward model reduces such degradation by roughly half for the Chat Hard subset in RewardBench.
Furthermore, when used in alignment, our robust reward models demonstrate better utility and lead to higher-quality outputs, winning in up to 59\% of instances against a standardly trained RM.
